% Avsnitt 6.1-6.7, ur lectures/L06/appendix/a_counters_and_shift_registers.md (A.1-A.7).

\section{Från register till räknare}
\label{c6:sec:counter}

Ett register (\chapref{c3:ch}) är en bank av D-vippor som lagrar vilket värde som än matas in vid
nästa stigande flank. En \textbf{räknare} är samma idé med en vridning: i stället för ett nytt värde
utifrån matar du registret med dess \emph{eget nuvarande värde plus ett}. Vid varje flank fångar det
\code{current_value + 1}, så det lagrade talet ökar en gång per klockcykel utan någon vidare
inmatning.

\subsection*{Bygg den för hand först}
Som med varje krets i den här boken: bygg den innan du skriver den. En 4-bitars räknare är tre
element, och att se den köra är hela poängen:
\begin{itemize}
  \item Ett \textbf{4-bitars register}: fyra D-vippor som delar en klocka, exakt registret från
    \secref{c3:sec:register}.
  \item En \textbf{adderare} och en \textbf{konstant \code{1}}: registrets fyra utgångar till
    adderarens ena ingång, konstanten till den andra.
  \item Adderarens summa kopplad tillbaka till registrets fyra \code{D}-ingångar. Den
    återkopplingen är hela kretsen: registret håller ett tal, adderaren beräknar det talet plus
    ett, nästa flank lagrar det. Det är \secref{c3:sec:seq}:s återkopplingsslinga igen, med en
    adderare i vägen i stället för en ren ledning.
\end{itemize}

\lecfig[0.7]{lectures/L06/appendix/images/counter_circuit.png}{En 4-bitars räknare: ett register
vars egen utgång, plus ett, kopplas tillbaka till dess ingång.}{c6:fig:countercircuit}

Bekräfta sedan, en klockflank i taget, med en period på \code{1000 ms}:
\begin{itemize}
  \item Se utgångarna stega \code{0000, 0001, 0010, 0011, ...}, en räkning per stigande flank.
  \item \textbf{Låt den köra förbi \code{1111}.} Den återgår till \code{0000} av sig själv, och
    ingenting i din krets säger åt den att göra det: det finns ingen jämförelse, ingen
    nollställningslogik, inget ”om räkningen står på maxvärdet” någonstans. Adderaren producerade
    \code{10000}, den femte biten hade ingenstans att ta vägen, och de fyra bitar som överlevde är
    \code{0000}. Det är \textbf{spill}, och i en räknare är det inte en bugg att arbeta runt utan
    den egenskap som resten av det här kapitlet och hela \chapref{c7:ch} vilar på.
  \item \textbf{Håll ögonen på \code{carry_out} på adderaren medan det sker.} Det är den femte
    biten, och ritningen visar exakt var den hamnar: hög under den enda flank där räkningen slår
    över, kopplad till ingenting, ut ur adderaren och stopp där. En bredare räknare skulle mata in
    den i ett femte steg; här kastas den bort, och att kasta bort den \emph{är} överslaget. Så
    ”räkningen slår över för att bitarna tar slut” är inte ett bildligt uttryck. Du kan peka på
    biten som tog slut.
\end{itemize}

Spara projektet. \Chapref{c7:ch} utgår från exakt den här kretsen och lägger till två saker till
den, och att öppna den igen är bättre än att rita om den.

\subsection*{Samma räknare i VHDL}
En signal och en process, precis som \chapref{c3:ch}:s register, bortsett från att högerledet
refererar till signalen själv:

\begin{vhdlcode}
signal counter: natural range 0 to 15;
...
process(clock, reset_s2_n) is
begin
    if (reset_s2_n = '0') then
        counter <= 0;
    elsif (rising_edge(clock)) then
        counter <= counter + 1;
    end if;
end process;
\end{vhdlcode}

\code{natural} är en fördefinierad heltalssubtyp i VHDL: ett icke-negativt heltal, till skillnad
från \code{std_logic}-typerna som modellerar bitar direkt. VHDL:s aritmetiska operatorer fungerar på
den så som vanlig heltalsmatematik gör, vilket är varför \code{counter + 1} inte behöver något
extra maskineri. Syntesen gör ändå om den till ett register av vippor, precis som ett
\code{std_logic_vector} skulle bli; skillnaden ligger enbart i hur du skriver och resonerar om
värdet.

Intervallbegränsningen \code{0 to 15} är inte bara dokumentation. Den talar om för syntesverktyget
att den här räknaren behöver fyra bitar, så att räkna förbi \code{15} faller tillbaka till \code{0}
gratis, enbart som en konsekvens av att bitarna tar slut, utan någon uttrycklig ”om counter = 15”-\!
logik vare sig i ritningen eller i VHDL-koden.

Det gratis överslaget har ett villkor värt att säga ut: det inträffar bara när den övre gränsen är
$2^N - 1$. \code{0 to 15} slår över gratis; \code{0 to 9} gör det inte, eftersom \code{10} inte är
en tvåpotensgräns, och en modulo-10-räknare behöver en uttrycklig jämförelse. Det är precis vad
Övning~\ref{c6:ex:counter} ber dig skriva.

\begin{warning}
\textbf{Ett ställe där simulatorn är oense med hårdvaran.} Det gratis överslaget är en egenskap hos
den \emph{syntetiserade} kretsen, och ritningen visar att det är verkligt. En
VHDL-\emph{simulator} är strängare än en ledning:
\begin{itemize}
  \item \code{natural range 0 to 15} intervallkontrolleras, så GHDL behandlar $15 + 1 = 16$ som ett
    fel och avbryter i stället för att slå över.
  \item Den syntetiserade hårdvaran har ingen sådan föreställning. Det finns inget sextonde värde
    för bitarna att anta, så de landar på \code{0000}, precis som dina fyra lysdioder gjorde.
\end{itemize}
Grindnätet och chippet slår båda över; bara simulatorn invänder, eftersom ett intervallbegränsat
heltal är ett löfte du gett och \code{16} bryter det. Regeln som följer är inte ”lita på
hårdvaran”: den är att allt du tänker verifiera bör säga vad det menar. Skriv ut jämförelsen, så
som \code{timer} och \code{serial_rx8} gör, så är ritningen, simulatorn och chippet alla överens.
\end{warning}

Det här överslaget som följer av bitbredden är räknarens enskilt viktigaste egenskap, och det är av
det \chapref{c7:ch} bygger en timer: en räknare som håller utkik efter ett bestämt räknarvärde och
höjer en flagga när det kommer.

% ------------------------------------------------------------------------------------------
\section{Skiftregister}
\label{c6:sec:shift}

Ett \textbf{skiftregister} är, återigen, N D-vippor, kopplade så att varje vippas utgång matar
\emph{nästa} vippas ingång i stället för sin egen. Vid varje klockflank flyttar varje bit en
position nedåt i kedjan, samtidigt:

\begin{textdrawing}
serial_in -> [D Q]-->[D Q]-->[D Q]-->[D Q] -> serial_out
              FF0     FF1     FF2     FF3
             bit 0   bit 1   bit 2   bit 3
              ^clock  ^clock  ^clock  ^clock   (alla delar samma klocka)
\end{textdrawing}

Efter en flank ligger \code{FF0}:s innehåll i \code{FF1}, \code{FF1}:s i \code{FF2}, och så vidare;
en ny bit kommer in i \code{FF0} från \code{serial_in}, och det som låg i den sista vippan faller av
änden som \code{serial_out}. Kör den i N flanker så har N seriella bitar vandrat hela vägen igenom.

\textbf{Ett håll, som används genomgående i den här boken.} Ett skiftregister kan flytta bitar åt
båda hållen, och båda är lika giltiga; det är att blanda dem inom en och samma design som skapar
förvirring. Så varje skiftregister här följer diagrammet ovan:
\begin{itemize}
  \item kedjan avbildas på ett \code{std_logic_vector} med \code{FF0} som \textbf{bit 0} och
    \code{FF(N-1)} som bit \code{N-1}.
  \item seriella data \textbf{kommer in vid bit 0} och \textbf{lämnar från bit \code{N-1}}.
  \item varje bit flyttar en position mot den \textbf{mest signifikanta} biten per flank.
\end{itemize}

I VHDL är det hållet ett enda uttryck, identiskt i varje exempel nedan:

\begin{vhdlcode}
shift_reg <= shift_reg(N-2 downto 0) & serial_in;
\end{vhdlcode}

\code{shift_reg(N-2 downto 0)} är allt utom den översta biten, flyttat upp en position; den översta
biten har ingenstans att ta vägen och faller av som \code{serial_out}; \code{serial_in} fyller den
lediggjorda bit 0.

När du möter ett skiftregister utanför den här boken: kontrollera dess håll innan du antar något.
Spegelbilden, \code{serial_in & shift_reg(N-1 downto 1)}, är precis lika vanlig och ser vid en
snabb blick nästan likadan ut.

Skiftregister namnges efter hur data kommer in och lämnar:

\begin{booktable}{@{}p{7.6em}p{6.6em}p{6.6em}L@{}}
  \thead{Utförande} & \thead{Data in} & \thead{Data ut} & \thead{Typisk användning} \\
  \midrule
  SISO & en bit per klocka & en bit per klocka & en ren fördröjningslinje \\
  \textbf{SIPO} & en bit per klocka & alla N bitar på en gång & att ta emot en ström av seriella
    bitar, t.ex. en inkommande byte på en enda ledning \\
  \textbf{PISO} & alla N bitar på en gång (laddade) & en bit per klocka & att sända ut N bitar över
    en enda ledning \\
  PIPO & alla N bitar på en gång & alla N bitar på en gång & i praktiken ett vanligt register
    (\secref{c3:sec:register}) \\
\end{booktable}

Det här kapitlet tar upp de två praktiskt användbara, SIPO och PISO.

% ------------------------------------------------------------------------------------------
\section{Serie in, parallell ut (SIPO)}
\label{c6:sec:sipo}

Ett SIPO-register är det du griper efter när bitar anländer en i taget på en enda ledning och du
behöver det ackumulerade värdet på en gång. Klassikern är att deserialisera en inkommande byte,
vilket är vad mottagarvägen i vilket seriellt protokoll som helst i grunden gör.

Vilken ände av byten som anländer först är varje protokolls eget beslut. SPI och CAN sänder som
förval den mest signifikanta biten först, vilket är vad \secref{c6:sec:shift}:s håll förutsätter:
den första biten som anländer hamnar i MSB när alla N har skiftats in. (Bitordningen i SPI går att
ställa om på de flesta styrkretsar, så kontrollera komponenten.) Ett protokoll som sänder den minst
signifikanta biten först, som UART, skiftar åt det spegelvända hållet, och dess idiom är också
spegelbilden.

Idiomet är en rad inuti en klockad process, \secref{c6:sec:shift}:s skiftuttryck ordagrant:

\begin{vhdlcode}
signal shift_reg: std_logic_vector(7 downto 0);
...
process(clock, reset_s2_n) is
begin
    if (reset_s2_n = '0') then
        shift_reg <= (others => '0');
    elsif (rising_edge(clock)) then
        shift_reg <= shift_reg(6 downto 0) & serial_in;
    end if;
end process;
parallel_out <= shift_reg;
\end{vhdlcode}

Efter 8 flanker har de mottagna bitarna helt ersatt registrets ursprungliga innehåll, och
\code{parallel_out} exponerar alla 8 på en gång.

% ------------------------------------------------------------------------------------------
\section{Parallell in, serie ut (PISO)}
\label{c6:sec:piso}

Ett PISO-register vänder på SIPO:s uppgift: ladda N bitar på en gång, sänd sedan ut dem en i taget.
Det är det du griper efter när ett fast mönster ska ut över en enda ledning, till exempel till en
kedja av lysdioder eller en display driven av ett externt skiftregisterchip.

Det behöver en styrsignal mer än SIPO: något som skiljer att \emph{ladda} ett nytt parallellt värde
från att \emph{skifta} ut det befintliga.

\begin{vhdlcode}
process(clock, reset_s2_n) is
begin
    if (reset_s2_n = '0') then
        shift_reg <= (others => '0');
    elsif (rising_edge(clock)) then
        if (load = '1') then
            shift_reg <= parallel_in;                          -- Parallel load.
        elsif (shift = '1') then
            shift_reg <= shift_reg(N-2 downto 0) & serial_in;  -- Shift toward the MSB.
        end if;
    end if;
end process;
serial_out <= shift_reg(N-1);
\end{vhdlcode}

Bit \code{N-1} är alltid nästa bit att lämna, så ett laddat värde går ut med den mest signifikanta
biten först: ladda \code{"1010"} i ett 4-bitars register så presenterar \code{serial_out} \code{1},
\code{0}, \code{1}, \code{0} under de följande cyklerna.

Lägg märke till hur lite som skiljer det här från \secref{c6:sec:sipo}:s SIPO.
\textbf{Skiftuttrycket är identiskt}, samma håll, samma rad VHDL. Bara två saker skiljer: PISO
lägger till en \code{load}-gren, och den tappar av \code{shift_reg(N-1)} som seriell utgång där
SIPO exponerar hela vektorn. ”SIPO” och ”PISO” namnger hur du \emph{kopplar} ett skiftregister,
inte två sorters hårdvara, vilket är varför \secref{c6:sec:shift} kunde beskriva skiftningen en gång
innan något av namnen dök upp.

Om en design behöver ett verkligt nytt \code{serial_in} vid varje skiftning, eller nöjer sig med att
återcirkulera det som faller av änden, avgörs av hur den är kopplad. Den återcirkulationen är precis
det trick som exemplet med den vandrande lysdioden i \chapref{c7:ch} använder, när det väl har en
timer som taktar skiftningen.

% ------------------------------------------------------------------------------------------
\section{Genomarbetat exempel: en 8-bitars seriemottagare}
\label{c6:sec:rx8}

Kapitlets två halvor möts i en modul. Ett skiftregister kan ta emot seriella bitar i all evighet men
kan inte tala om \emph{när} en komplett byte anlänt, eftersom det inte har någon föreställning om
”hur många bitar hittills”. Det är ett räkneproblem, och \secref{c6:sec:counter} löste det redan.
\code{serial_rx8} kombinerar de två, plus en \code{data_ready}-puls på en cykel när den åttonde
biten landar.

Det här är ingen leksak: det är främre halvan av varje seriemottagare som finns. Mottagarvägarna i
SPI, CAN och UART börjar alla med ett skiftregister och en biträknare, och skiljer sig bara i vad
som avgör när en bit är giltig.

\begin{booktable}{@{}p{6.6em}p{3.4em}p{9em}L@{}}
  \thead{Port} & \thead{Riktn.} & \thead{Typ} & \thead{Betydelse} \\
  \midrule
  \code{clock} & in & \code{std_logic} & Systemklocka. \\
  \code{reset_s2_n} & in & \code{std_logic} & Aktiv låg, \textbf{redan synkroniserad} reset;
    nollställer registret och biträkningen. Asynkron, så den nollställer dem utan att någon
    klockflank är inblandad. \\
  \code{shift_enable} & in & \code{std_logic} & Hög under exakt en klockcykel per inkommande bit. \\
  \code{serial_in} & in & \code{std_logic} & Den inkommande biten, samplad vid varje aktiverad
    stigande flank. \\
  \code{data_out} & out & \code{std_logic_vector(7 downto 0)} & Byten som mottagits hittills, mest
    signifikanta biten först. \\
  \code{data_ready} & out & \code{std_logic} & Puls på en cykel när den åttonde biten landar. \\
\end{booktable}

\begin{aside}
\textbf{Resetporten är den synkroniserade, och \code{_s2} säger det.} Den heter \code{reset_s2_n}
snarare än \code{reset_n} eftersom det som hör hemma på den är utgången från en \code{reset_sync}
(\secref{c4:sec:resetsync}), modulen du skrev i Övning~\ref{c4:ex:resetsync}: aktivera asynkront,
släpp synkront. \textbf{Instansiera en i den design som omsluter den här mottagaren och koppla dess
\code{reset_s2_n} hit}, vilket är precis vad \secref{c6:sec:shiftdemo}:s omslagsmodul på kortet
gör. Koppla i stället en tryckknapp rakt in på den här porten så får du den kapplöpning
\chapref{c4:ch} finns till för att förhindra: skiftregistret och biträknaren är ett dussin vippor,
och en asynkront \emph{släppt} reset låter dem lämna reset vid olika flanker, så mottagaren kan
börja räkna en byte en cykel innan den börjar skifta in en.

\code{counter}, \code{sipo8} och \code{piso8} i övningarna tar \code{reset_s2_n} av samma skäl.
Ingen av dem är en design du sätter på ett kort för sig själv: en räknare, ett skiftregister och en
mottagare är delar, och en del har rätt att anta att designen omkring den har skött
synkroniseringen. Så hela det här kapitlet tar den synkroniserade reset-signalen, och bara
omslagsmodulerna i \secref{c6:sec:counterdemo} och \ref{c6:sec:shiftdemo}, som äger pinnarna, får
någonsin se en rå \code{reset_n}.

En \textbf{toppnivå} namnges tvärtom: den tar \code{reset_n} och instansierar \code{reset_sync}
själv, eftersom det är den som håller pinnen. \code{blinker} och \code{walking_led} i
\chapref{c7:ch}, och båda capstones i \chapref{c8:ch}, är skrivna så. Skaffa dig vanan att läsa
\code{_s2} som ett krav på anroparen, för härifrån och framåt är det mesta du skriver en
subkomponent.
\end{aside}

\vhdlfile{lectures/L06/serial_rx8/serial_rx8.vhd}

Vid varje stigande flank där \code{shift_enable = '1'} skiftas \code{serial_in} in i bit 0 med
\secref{c6:sec:sipo}:s idiom oförändrat, så den första biten som anländer har vandrat upp till bit 7
när den åttonde landar, och en intern biträknare av typen \code{natural range 0 to 7} räknar upp.
Vid den flank där den räknaren redan står på \code{7} pulsar \code{data_ready} hög under en cykel
och räknaren nollställs tillbaka till \code{0}, redo för nästa byte. \code{data_out} drivs av en
konkurrent tilldelning från skiftregistret, så den kompletta byten ligger på den under den cykel då
\code{data_ready} är hög.

Lägg märke till att räknaren \emph{nollställs} här snarare än lämnas att slå över: \code{0 to 7} är
ett tvåpotensintervall så den skulle slå över gratis, men att skriva ut det håller simulatorn
överens med hårdvaran, enligt \secref{c6:sec:counter}.

Två detaljer är värda att dröja vid, eftersom båda är mönster du mött mer än en gång:
\begin{itemize}
  \item \textbf{\code{data_ready} är en puls, inte en nivå.} Den drivs låg villkorslöst högst upp i
    den klockade grenen och höjs bara vid den flank som fullbordar en byte: samma encykelsform som
    flankdetektorn i \secref{c3:sec:edge}. En konsument som missar den har missat byten, vilket är
    precis därför den är parad med ett \code{data_out} som \emph{håller} sitt värde efteråt. En
    riktig kringkrets skulle sätta ett hållregister eller en liten FIFO här så att en långsam
    konsument inte kan förlora data.
  \item \textbf{\code{shift_enable} grindar biträknaren, inte bara skiftningen.} Båda går framåt
    tillsammans eller ingendera, så ett avbrott i den inkommande strömmen pausar mottagaren mitt i
    en byte i stället för att korrumpera den, och den återupptar på exakt den bit den väntade på.
\end{itemize}

\paragraph{Att kontrollera den}
Till exemplet hör en utdelad \file{serial_rx8_tb.vhd}, som skiftar in två byte med en paus emellan
och kontrollerar att \code{data_ready} pulsar en gång per byte, aldrig för tidigt:

\begin{shell}
cd lectures/L06/serial_rx8
ghdl -a --std=93 serial_rx8.vhd serial_rx8_tb.vhd
ghdl -e --std=93 serial_rx8_tb
ghdl -r --std=93 serial_rx8_tb --assert-level=error --stop-time=10ms
\end{shell}

Taktad på \code{50 MHz} och kontinuerligt aktiverad sväljer den här mottagaren en byte på 160~ns,
vilket ingen lysdiod kan visa dig. Det är därför testbänken spelar roll här på ett sätt den inte
gjorde för en blinkande lysdiod, och så är det för det mesta av verklig digitalkonstruktion. Men det
är inget skäl att hoppa över kortet, för \code{shift_enable} är just utvägen: se
\secref{c6:sec:shiftdemo}.

% ------------------------------------------------------------------------------------------
\section{På kortet: att se en räknare slå över}
\label{c6:sec:counterdemo}

Räknaren demonstreras på DE0-CV under passet, byggd på samma sätt som varje design sedan
\chapref{c1:ch}: ett Quartus Prime Lite-projekt mot \code{5CEBA4F23C7N}, \file{.vhd}-filen tillagd,
kompilera, tilldela varje port en fysisk pinne i Pin Planner, kompilera sedan om och programmera
kortet.

En sak måste ändras först, och det är värt att se varför. Taktad på \code{50 MHz} går en 4-bitars
räknare igenom alla sexton tillstånd och slår över var \code{320 ns}, ungefär tre miljoner gånger i
sekunden. En lysdiod driven direkt från den blinkar inte, den lyser helt enkelt med halv styrka. För
att se ett överslag måste du titta på en mycket långsammare bit: räkna i ett bredare register, säg
28 bitar, och tilldela de \textbf{översta} fyra bitarna till lysdioderna.

Demonstrationen är därför en omslagsmodul kring två moduler snarare än övningens entitet för sig:
en \code{reset_sync} som tar den råa \code{reset_n} från pinnen, och den breddade räknaren som tar
sin \code{reset_s2_n} därifrån. Räknarens port heter \code{reset_s2_n} just för att den här
kopplingen ska vara den enda som typkontrollerar, och \secref{c6:sec:shiftdemo} upprepar mönstret
med en tredje modul.

\begin{itemize}
  \item \code{clock} till 50 MHz-oscillatorn, \code{reset_n} till en tryckknapp (in i
    \code{reset_sync}, aldrig rakt in i räknaren), \code{count(27 downto 24)} till fyra lysdioder.
  \item Bit 24 ändras en gång var $2^{24}$ klockcykler, ungefär en tredjedels sekund, så de fyra
    lysdioderna räknar synligt från \code{0000} upp till \code{1111} och sedan, utan att något i
    kretsen gör det, tillbaka till \code{0000}. Ett helt varv tar ungefär fem sekunder.
\end{itemize}

Att dela en klocka genom att titta på en hög bit är den grövsta tänkbara timern, och dess
begränsning är värd att namnge redan nu: perioden är låst till en tvåpotens, och det enda sättet
att ändra den är att välja en annan bit. \Chapref{c7:ch} ersätter den med en räknare jämförd mot
ett målvärde du väljer, vilket är allt en timer egentligen är.

% ------------------------------------------------------------------------------------------
\section{På kortet: att skifta in en byte för hand}
\label{c6:sec:shiftdemo}

\Secref{c6:sec:rx8} sa att en byte anländer för fort för att kunna följas. Lösningen är inte att
sakta ned klockan, vilket vore mjukvaruinstinkten och är fel drag i hårdvara: klockan är den enda
sak i en synkron design som du inte pillar på. \code{serial_rx8} har redan rätt styrning, och det är
\code{shift_enable}, som ”grindar biträknaren, inte bara skiftningen”. Håll den låg så stannar
mottagaren tvärt, mitt i en byte, hur länge du vill. Höj den under exakt en cykel så går exakt en
bit in.

Driv alltså \code{shift_enable} från en \textbf{knapp} i stället för att hålla den hög. Ett tryck,
en bit.

Demonstrationen är en liten omslagsmodul kring tre moduler du redan har, och ingen ny logik utöver
en enda vippa:
\begin{itemize}
  \item \code{reset_sync} (\secref{c4:sec:resetsync}), som gör den råa \code{reset_n} till
    \code{reset_s2_n} för allt annat.
  \item \code{button_sync} (\secref{c4:sec:buttonsync}) med \code{COUNT = 1}, som gör ett tryck
    till den encykelspuls som \code{shift_enable} vill ha. Det här är ingen valfri utsmyckning: en
    rå knapp kopplad till \code{shift_enable} skulle skifta in en slumpmässig handfull bitar per
    tryck, en per klockcykel så länge kontakterna studsar.
  \item \code{serial_rx8} själv, oförändrad, med sin \code{reset_s2_n} tagen från
    \code{reset_sync} ovan snarare än från knappen. Det är hela skälet till att porten heter
    \code{reset_s2_n}: omslagsmodulen äger den råa \code{reset_n}, och mottagaren får bara någonsin
    se den synkroniserade.
  \item En vippa som låser \code{data_ready}, satt när pulsen anländer och nollställd av reset.
    Utan den finns ingenting att se: \code{data_ready} är hög i 20~ns, och det kan ingen lysdiod
    visa dig.
\end{itemize}

På kortet: \code{clock} till 50 MHz-oscillatorn, \code{reset_n} till en tryckknapp,
\code{serial_in} till en strömbrytare (det är den bit du är på väg att skicka), skiftknappen till
en andra tryckknapp, \code{data_out(7 downto 0)} till åtta lysdioder, och den låsta
\code{data_ready} till en nionde.

Ställ nu strömbrytaren, tryck på knappen, och se en bit komma in vid lysdioden längst till höger
och hela mönstret stega åt vänster. Skicka \code{10110010} en bit i taget så byggs byten upp över
lysdioderna framför dig, mest signifikanta biten först, precis som \secref{c6:sec:sipo}:s idiom
säger att den måste. Vid det åttonde trycket tänds den nionde lysdioden och biträknaren börjar om.

Tre saker som det här gör synliga och som testbänken bara kan hävda:
\begin{itemize}
  \item \textbf{Bitordningen.} Den första biten du skickar hamnar i bit 7. Att läsa av det ur en
    vågform är ett tjatgöra; att se den vandra över en rad lysdioder är det inte.
  \item \textbf{Vad \code{shift_enable} är till för.} Släpp knappen en minut mitt i en byte så rör
    sig ingenting, ingenting går förlorat, och den åttonde biten fullbordar ändå byten. Det är
    egenskapen ”pausar mitt i en byte i stället för att korrumpera den”, demonstrerad snarare än
    påstådd.
  \item \textbf{Puls mot nivå, en gång till.} Skälet till att en vippa måste läggas till för
    \code{data_ready} är hela den distinktionen, och det är samma skäl som flankdetektorn i
    \secref{c3:sec:edge} fanns till för.
\end{itemize}
